\documentclass[twocolumn]{aastex631}
% Auto-generated by jansky_research.atlas3i.survey_macros — do not edit by hand.
\newcommand{\aiTotNodes}{60}
\newcommand{\aiTotScans}{360}
\newcommand{\aiTotBands}{4}
\newcommand{\aiTotRawHits}{1\,938\,860}
\newcommand{\aiTotSurvivors}{294}
\newcommand{\aiTotConfirmed}{0}
\newcommand{\aiSatCoherent}{9}
\newcommand{\aiEirpMw}{99.2}
\newcommand{\aiDistanceAu}{1.798}
\newcommand{\aiLNodes}{6}
\newcommand{\aiLBandLo}{939}
\newcommand{\aiLBandHi}{2064}
\newcommand{\aiLRawHits}{1\,124\,011}
\newcommand{\aiLSurvivors}{261}
\newcommand{\aiLConfirmed}{0}
\newcommand{\aiSNodes}{6}
\newcommand{\aiSBandLo}{1652}
\newcommand{\aiSBandHi}{2776}
\newcommand{\aiSRawHits}{354\,423}
\newcommand{\aiSSurvivors}{29}
\newcommand{\aiSConfirmed}{0}
\newcommand{\aiCNodes}{23}
\newcommand{\aiCBandLo}{3939}
\newcommand{\aiCBandHi}{8252}
\newcommand{\aiCRawHits}{15\,735}
\newcommand{\aiCSurvivors}{0}
\newcommand{\aiCConfirmed}{0}
\newcommand{\aiXNodes}{25}
\newcommand{\aiXBandLo}{7689}
\newcommand{\aiXBandHi}{12376}
\newcommand{\aiXRawHits}{444\,691}
\newcommand{\aiXSurvivors}{4}
\newcommand{\aiXConfirmed}{0}

\newcommand{\code}[1]{\texttt{#1}}

\shorttitle{Reproducing the BL 3I/ATLAS nondetection}
\shortauthors{Barbere}
\begin{document}

\title{An Independent Reproduction of the Breakthrough Listen 3I/ATLAS Technosignature
Nondetection from the Public Green Bank Telescope Archive}
\author[0009-0008-3289-4447]{Joseph Barbere}
\affiliation{Independent researcher}

\begin{abstract}
Breakthrough Listen (BL) searched the interstellar object 3I/ATLAS for radio
technosignatures with the Green Bank Telescope at 1--12~GHz during its 2025 December
close approach and reported a nondetection down to $\sim$100~mW equivalent isotropic
radiated power (EIRP), releasing all data publicly. Technosignature nulls are almost
never independently re-examined. I reprocess the complete released archive ---
\aiTotNodes\ unique 187.5~MHz compute-node cadences across the four receivers --- with
a fully independent open-source pipeline: a physical-units brute-force de-Doppler search
($\pm4$~Hz\,s$^{-1}$, $16\sigma$ against robust per-drift noise), the ABACAD on/off sky
filter, a per-candidate drift-coherence vet on postage stamps re-read from the archive,
and an ITU satellite-allocation exclusion. The null reproduces in full: \aiTotRawHits\
raw hits reduce to \aiTotSurvivors\ two-position survivors and \aiTotConfirmed\
confirmed candidates, and the implied narrowband limit ($\le\aiEirpMw$~mW at L band for
$d=\aiDistanceAu$~au) is algebraically identical to the published one at unit
dedrifting efficiency. The survivors form a physically interpretable taxonomy of
filter evasion --- sub-threshold terrestrial carriers, chance coincidences in dense hit
lists, and satellite downlinks that defeat two-position geometry by construction ---
independently recovering the design rationale of BL's interference rejection. The
entire reproduction cost $\approx$95~h of active compute and $\approx$3.7~TB of
transfer on a single desktop computer.
\end{abstract}

\keywords{Technosignatures (2128) --- Search for extraterrestrial intelligence (2127)
--- Radio astronomy (1338) --- Interstellar objects (52)}

\section{Introduction}
Radio technosignature searches report, almost without exception, nondetections
\citep[e.g.][]{gajjar2021,choza2024}. Those nulls carry quantitative content --- EIRP
upper limits over a defined haystack volume --- yet they are essentially never
independently reproduced: the datasets are large, the pipelines are specialised, and a
null invites no scrutiny. The 2025 apparition of the interstellar object 3I/ATLAS
produced an unusually visible cluster of such results: BL observed it with the GBT
across 1--12~GHz at its closest approach and reported no candidate signals above
$\sim$100~mW EIRP \citep{jacobsonbell2025}, with independent-\emph{observation} nulls
from the Allen Telescope Array \citep{ata3i} and FAST \citep{fast3i}. All BL data were
released publicly at \url{https://bldata.berkeley.edu/ATLAS}.

This paper asks the complementary question: does the published null survive an
independent \emph{reanalysis} of the same released data with a pipeline sharing no code
with the original? Same-data reproduction cannot confirm the observation itself, but it
tests everything downstream of the telescope: the search's completeness against its
stated parameters, the interference-rejection logic, and the sensitivity arithmetic.
For technosignature claims --- where a future positive would demand extraordinary
scrutiny --- demonstrating that the null pipeline audit is practical seems worth doing
once, explicitly, on a high-profile target.

\section{Data}
BL observed 3I/ATLAS on MJD 61027 (2025 December 18) with four GBT receivers (L, S, C,
X) in consecutive 6$\times$5-min ABACAD on/off cadences, recording 187.5~MHz per
compute node \citep{jacobsonbell2025}. The released fine-frequency products
($\sim$2.79~Hz channels, $\sim$18.25~s integrations, $\sim$10~GB per scan per node)
were mapped by reading every node's header remotely (HTTP range requests): L spans
\aiLBandLo--\aiLBandHi~MHz over \aiLNodes\ nodes, S \aiSBandLo--\aiSBandHi~MHz
(\aiSNodes), C \aiCBandLo--\aiCBandHi~MHz (\aiCNodes), and X
\aiXBandLo--\aiXBandHi~MHz (\aiXNodes). Twelve additional C/X-band recordings duplicate
bank-boundary coverage byte-for-byte and were searched once. The recorded spans exceed
the passbands the original analysis quotes (L: 1.1--1.9~GHz; X: 7.6--11.7~GHz); the
band-edge roll-off regions are included here and flagged where they matter. In total:
\aiTotNodes\ unique node-cadences, 366 fine-resolution scans.

\section{Methods}
The pipeline (\code{atlas3i} in \code{jansky-research}; \citealt{janskyresearch}) is
pure NumPy and shares no code with \code{turboSETI}. Each scan is processed in bounded
frequency chunks: per-channel bandpass normalization by the time-median, excision of
channels deviating $>10$ robust~$\sigma$ from the band median (removing the rawspec
band-centre DC spike and persistent carriers), then a brute-force shift-and-sum
de-Doppler over 129 drift rates spanning $\pm4$~Hz\,s$^{-1}$ in physical units.
Detections are peaks above $16\sigma$ of the per-drift-row median/MAD noise, with
neighbour suppression. Candidates must then survive three physically distinct filters:

\emph{(i) ABACAD geometry.} A hit must appear, extrapolated along its own drift rate,
in every on-target scan and no off-target scan.

\emph{(ii) Drift coherence.} For each survivor the pipeline re-reads $\pm$2048-channel
postage stamps from all six scans and demands the tone track its fitted drift through
the full $\sim$27-min cadence, with direct off-scan cleanliness at $8\sigma$ ---
subsuming the original search's asymmetric ($16\sigma$/$10\sigma$) on/off thresholds,
which this reproduction otherwise does not replicate (a stricter symmetric choice that
passes \emph{more} borderline interference forward; conservative for a null).
Exactly-zero drift is flagged terrestrial.

\emph{(iii) Frequency allocation.} Survivors inside ITU satellite-downlink allocations
(GNSS bands, Iridium, Inmarsat/MSS, S-DARS, Ku FSS/DBS) cannot be confirmed: satellite
transmitters are genuinely in the sky and intermittent, so no two-position filter
rejects them by geometry --- the same rationale as BL's exclusion masks.

Every stage is validated offline on a synthetic cadence containing an injected drifting
tone (recovered), an always-on interferer (rejected by geometry), and a DC spike
(excised); the detector's false-positive rate and recovery completeness were calibrated
in a companion injection-recovery benchmark \citep{janskyresearch}. The 3I/ATLAS
geocentric distance at the cadence epoch, $d=\aiDistanceAu$~au, is taken from JPL
Horizons.

\section{Results}
\begin{deluxetable}{lcccc}
\tablecaption{Full-survey vetting funnel by receiver band.\label{tab:funnel}}
\tablehead{\colhead{band} & \colhead{nodes} & \colhead{raw hits} &
\colhead{on/off survivors} & \colhead{confirmed}}
\startdata
L & \aiLNodes & \aiLRawHits & \aiLSurvivors & \aiLConfirmed \\
S & \aiSNodes & \aiSRawHits & \aiSSurvivors & \aiSConfirmed \\
C & \aiCNodes & \aiCRawHits & \aiCSurvivors & \aiCConfirmed \\
X & \aiXNodes & \aiXRawHits & \aiXSurvivors & \aiXConfirmed \\
\hline
total & \aiTotNodes & \aiTotRawHits & \aiTotSurvivors & \aiTotConfirmed \\
\enddata
\end{deluxetable}

The null reproduces across the full band (Table~\ref{tab:funnel},
Figure~\ref{fig:survey}): no candidate survives all three filters anywhere in
1--12~GHz. The hit-density contrast is extreme --- from $\sim$2$\times10^{5}$ hits per scan in
the worst L-band node down to individual C-band scans recording zero raw hits ---
and, separately, the C band produced \emph{zero} two-position survivors across its
4.3~GHz. The narrowband sensitivity,
$S_\mathrm{min} = 16\,\mathrm{SEFD}/\sqrt{n_\mathrm{pol}\,\delta\nu\,\tau}$ and
$\mathrm{EIRP} = 4\pi d^{2} S_\mathrm{min}\,\delta\nu$, gives
$\le\aiEirpMw$~mW at L band --- and reduces algebraically to the published formulation
\citep{gajjar2021} at unit dedrifting efficiency $\beta$, so the agreement with the
paper's $\sim$100~mW is exact rather than coincidental (our limit is marginally
optimistic at the highest drift rates, where $\beta<1$).

\subsection{The filter-evasion taxonomy}
The \aiTotSurvivors\ two-position survivors decompose cleanly. \emph{Sub-threshold
terrestrial carriers} (L band, 1302--1308~MHz, aeronautical radionavigation) passed the
on/off filter by sitting below the search threshold in off scans while above it in on
scans; the stamp vet detects them directly in the off data. \emph{Chance coincidences}
arise in the densest hit lists ($\sim$10$^{5}$ per scan) at the $\pm$32-channel matching
tolerance; none track their own drift. \emph{Satellite downlinks} are the instructive
class: \aiSatCoherent\ candidates passed \emph{both} geometric filters --- all in the
Iridium (1616--1626.5~MHz) or Inmarsat (1525--1559~MHz) bands with LEO-Doppler-scale
drifts and duplicate frequencies at multiple drift rates. X band contributed only four
survivors, all zero-drift Ku-band satellite-television tones above the analysed
passband. That an independent pipeline, built without reference to BL's interference
strategy, is forced to rediscover the same three-axis rejection design (geometry,
coherence, allocation) is itself a validation of the original search's architecture.

\subsection{Computational cost}
% The figures in this subsection are run-log measurements (survey/atlas3i-findings.md),
% not macro-generated from the result JSONs — update them there if the runs are redone.
Measured from run logs: $\approx$95~h of active compute (L $\approx$2.2~h per node;
S/C/X $\approx$1.5~h) and $\approx$3.7~TB transferred, with peak disk of one cadence
($\approx$60~GB) under a fetch--search--vet--delete cycle, on one desktop (6-core Ryzen
5 5600X, 62~GB RAM), CPU only. Auditing a published technosignature search from its
public archive requires no institutional resources.

\section{Discussion}
Three caveats bound the claim. First, same-data reproduction validates the analysis,
not the observation: a signal absent from the recorded voltages is beyond any
reanalysis. Second, detector conventions differ --- the per-drift-row MAD threshold
used here fires far more often than \code{turboSETI}'s normalization in dense
interference (our L-band raw hit count alone exceeds the original survey's full-band
total), so raw hit counts are not comparable across pipelines; for a null this is the
conservative direction, admitting more false positives into vetting and still
confirming none. Third, the recorded band edges searched here (e.g.\ 939~MHz,
12.4~GHz) lie in receiver roll-off the original team excluded; all survivors unique to
those regions failed vetting regardless.

With those bounds, the conclusion is direct: the BL 3I/ATLAS GBT nondetection is
robust to full independent reanalysis of the released data across the complete
1--12~GHz coverage, at matching quoted sensitivity (at L band) --- and the audit is cheap enough
to become routine for high-visibility technosignature results.

\software{\code{jansky-research} \citep{janskyresearch}, \code{jansky} \citep{jansky},
NumPy, h5py, astropy. Analysis assisted by Claude \citep{anthropicclaude}; all results
verified by the pipeline's test suite.}

\bibliography{refs}
\bibliographystyle{aasjournal}

\end{document}
